\chapter{Produit scalaire et orthogonalité}
\section*{Exercice 16.1 -- Coordonnées et angle}
$u\cdot v=3-2=1$, $\|u\|=\sqrt5$, $\|v\|=\sqrt{10}$. Le cosinus vaut $1/\sqrt{50}>0$ ; l'angle non orienté est donc aigu. Il n'est pas nécessaire de l'approcher en degrés pour le qualifier.

\section*{Exercice 16.2 -- Orthogonalité}
Le produit scalaire vaut $3t-12$. Il s'annule si et seulement si $t=4$. Les deux vecteurs sont alors non nuls et leur angle est droit.

\section*{Exercice 16.3 -- Identité de norme}
Développer les deux carrés par bilinéarité donne respectivement $\|u\|^2+2u\cdot v+\|v\|^2$ et $\|u\|^2-2u\cdot v+\|v\|^2$. Leur somme annule les termes croisés. Cette identité du parallélogramme reste vraie dans tout espace à produit scalaire.

\section*{Exercice 16.4 -- Projection signée}
Le produit vaut $4\times3\times(-1/2)=-6$. La projection algébrique vaut $(u\cdot v)/\|u\|=-3/2$. Son signe négatif indique que le projeté est du côté opposé au sens de $u$.

\section*{Exercice 16.5 -- Cauchy--Schwarz}
Dans $\R^3$, appliquer Cauchy--Schwarz aux vecteurs $(x,y,z)$ et $(1,1,1)$. Leur produit vaut $x+y+z$ et la seconde norme au carré vaut $3$. L'égalité a lieu si et seulement si les vecteurs sont colinéaires, donc $x=y=z$. La preuve par un carré du chapitre fonctionne aussi en dimension $3$.

\section*{Exercice 16.6 -- Base non orthonormée}
$\|e_1+e_2\|^2=1+1+2\cos(\pi/3)=3$. Les coordonnées dans cette base sont $(1,1)$, dont la somme des carrés vaut seulement $2$. La formule usuelle en coordonnées exige une base orthonormée ; ici le terme croisé est non nul.
